Для аналізу та тестування алгоритму бралися дані за червень 2019 року
для Київської області. Було проведено аналіз їхньої
відповідності та можливості подальшої обробки, оскільки якщо в одній
серії за місяць будуть знімки, де хмарність становить більше 50%, то на
виході алгоритму буде композит, в якого багато пікселів матимуть
чорний колір.

Тобто постала проблема відбору даних, що більш-менш підходять під
реалізацію методу. Також цей процес потрібно було повністю
автоматизувати, оскільки не зручно власноруч шукати потрібні знімки.

Цю проблему було вирішено за допомогою Open Data
Cube. ODC надав змогу автоматично та швидко завантажувати дані,
увівши параметри часу, широти, висоти та супутника.

Було написано процедури, які давали відсоток хмарності, що покращило
обробку, і таким чином відкидалися знімки за той час, де відсоток хмарності
перевищував 35%.
